\section{Phạm vi của phương pháp}\label{sec:02-scope}
Phương pháp trong chương này có ba giới hạn cần giữ khi đọc kết quả. Thứ nhất, dữ liệu phân cụm và xếp hạng là bán tổng hợp hoặc tổng hợp, nên kết quả RQ1--RQ3 mô tả hành vi của thuật toán dưới các điều kiện mô phỏng có kiểm soát, không phải hiệu quả trên báo cáo cứu hộ thực tế. Thứ hai, bộ ảnh do nhóm thu thập và gán nhãn, chưa có giấy phép cho từng ảnh, chưa có mẫu chụp thực địa độc lập và chưa đo độ đồng thuận giữa người gán nhãn. Thứ ba, các hằng số của $Q_i$ và điểm ưu tiên là heuristic của nhóm, chưa được chuyên gia cứu hộ xác nhận.

Đối với thông tin khẩn cấp, đề tài sử dụng Logistic Regression từ các lựa chọn có cấu trúc trên biểu mẫu, thay cho phương án phân tích văn bản tiếng Việt tự do dự kiến ban đầu. Các dấu hiệu được người dùng lựa chọn được mã hóa thành đầu vào để mô hình ước lượng mức độ khẩn cấp.

Nội dung nghiên cứu được đối chiếu với thuyết minh ở Mục~\ref{sec:proposal-comparison}; giới hạn của các phép thử được trình bày ở Mục~\ref{sec:04-evaluation-limitations}.
